% ---------------------------------------------------------------------------
% How the physical information of one choice class migrates between its
% subgraphs.  Companion text to figures/choice_migration_physics.svg.
%
% The equations follow the implementation in src/phys/place/physPlace.c:
%   Phys_PlaceClassAnchors   (anchor, eq. anchor)
%   Phys_PlaceSeedAlt        (seed grid, eqs. grid/seed)
%   Phys_PlaceEvalAlt        (local solve, eq. local)
%   Phys_PlaceClassCost      (cost, eq. cost)
%   Phys_PlaceSelectRounds   (switch, eq. switch)
%   Phys_PlaceCollectNets    (net rebuild, eq. migrate, last line)
% ---------------------------------------------------------------------------

\paragraph{Notation.}
A choice class~$c$ owns the set of alternatives~$A_c$; the alternative
$a\in A_c$ is a cone~$C_a\subseteq V$ of the bound netlist and
$n_a=|C_a|$ is its size.  $\sigma_c\in A_c$ is the selected alternative,
cell~$i$ is placed at $\mathbf p_i=(x_i,y_i)$ and occupies
$W_i\times H_i$, $w_s\times h_s$ is the site pitch, and
$\mathcal N(C_a)$ is the set of nets whose driver or sinks lie in the
cone.  $V_{\mathrm{on}}$ is the set of active nodes of the placement.

\paragraph{Anchor.}
The alternatives of a class share their fanins and their output, so the
only boundary at which they can be compared is the placement of that
output.  The anchor is the centroid of the cells of the selected cone,
\begin{equation}
  \boldsymbol\alpha_c
  \;=\; \frac{1}{n_{\sigma_c}}\sum_{i\in C_{\sigma_c}} \mathbf p_i ,
  \label{eq:anchor}
\end{equation}
and the value is stored on every member literal of the class, so a
candidate can be seeded around it and the next round can find it again.

\paragraph{Seed.}
A candidate is laid out on a grid of one site pitch per cell, centred on
the anchor,
\begin{equation}
  m_a=\lceil\sqrt{n_a}\rceil,\qquad
  q_i=i\bmod m_a,\qquad
  r_i=\left\lfloor i/m_a\right\rfloor,\qquad
  i=0,\dots,n_a-1,
  \label{eq:grid}
\end{equation}
\begin{equation}
  \mathbf p_i^{(0)} =
  \Big(\,
    \alpha_c^{x}+\big(q_i-\lfloor\tfrac{m_a-1}{2}\rfloor\big)w_s,\;
    \alpha_c^{y}+\big(r_i-\lfloor\tfrac{\lfloor (n_a-1)/m_a\rfloor}{2}\rfloor\big)h_s
  \,\Big).
  \label{eq:seed}
\end{equation}
The index is the position of the cell in the cone; an object of the cone
that is not a cell keeps its index and leaves a hole in the grid.  The
seeding never reads $W_i$: a wide cell is seeded like a narrow one, so
cells may overlap until the global re-placement legalises them.

\paragraph{Local solve.}
With the cone free and every other cell fixed, the placement solves the
same clique model as the global placement, restricted to the incident
nets, and adds a small anchor term to the seed so that a cell no net
constrains stays where the seeding put it,
\begin{equation}
  \min_{\{\mathbf p_i\}_{i\in C_a}}
  \sum_{q\in\mathcal N(C_a)} \frac{1}{|q|}
  \sum_{\{u,v\}\subseteq q}\|\mathbf p_u-\mathbf p_v\|^2
  \;+\;\varepsilon\sum_{i\in C_a}\|\mathbf p_i-\mathbf p_i^{(0)}\|^2 ,
  \qquad
  \varepsilon = 10^{-3}\Big(\frac{1}{n_a}\sum_{i\in C_a} d_i+1\Big),
  \label{eq:local}
\end{equation}
where $d_i$ is the diagonal of the clique Laplacian at cell~$i$ and a
fixed cell $v\notin C_a$ contributes its current position to the
right-hand side.

\paragraph{Cost and switch.}
The cost of an alternative is the exact half-perimeter wirelength of the
nets it takes part in, over the active pins,
\begin{equation}
  \kappa(a)=\sum_{q\in\mathcal N(C_a)}\Big[
    \max_{v\in q\cap V_{\mathrm{on}}} x_v-\min_{v\in q\cap V_{\mathrm{on}}} x_v
    +\max_{v\in q\cap V_{\mathrm{on}}} y_v-\min_{v\in q\cap V_{\mathrm{on}}} y_v
  \Big].
  \label{eq:cost}
\end{equation}
A candidate replaces the selection only when it beats the incumbent by
more than the tolerance, which keeps a tie from oscillating between two
alternatives,
\begin{equation}
  \sigma_c \leftarrow \arg\min_{a\in A_c}\kappa(a)
  \quad\text{iff}\quad
  \min_{a\in A_c}\kappa(a) < (1-\tau)\,\kappa(\sigma_c),
  \qquad \tau=0.01 .
  \label{eq:switch}
\end{equation}
The incumbent is measured at its globally placed coordinates and the
challenger after its local solve, so the comparison is between a global
configuration and a locally re-placed one.

\paragraph{What migrates.}
When $\sigma_c$ changes, the physical information of the class moves from
one subgraph to the other in a single step:
\begin{equation}
\begin{aligned}
  \mathbf p_i &\leftarrow \text{position after the local solve},
     && i\in C_{\sigma_c'},\\
  A_c &\leftarrow \sum_{i\in C_{\sigma_c'}} W_i\,H_i,\\
  \mathrm{ref}(v) &\leftarrow \sigma_c',
     && \text{for every cell $v$ outside the class}\\
  &&& \text{whose fanin referenced a member},\\
  \boldsymbol\alpha_c &\leftarrow \frac{1}{n_{\sigma_c'}}
     \sum_{i\in C_{\sigma_c'}}\mathbf p_i .
\end{aligned}
\label{eq:migrate}
\end{equation}
Only the selected cone takes part in the global problem: an inactive
alternative is neither moved nor counted as a pin of any net, so the
area, the density and the core follow $A_c$.  The fanout redirection is
what keeps the class output net attached to the new member.  After a round that
switched a class the net lists are rebuilt, the new representative graph
is placed globally from the coordinates of the round, and the anchors
are recomputed.

\paragraph{Equal and unequal sizes.}
If $n_a=n_b$ for all $a,b\in A_c$, the seeded footprint
$m\times\lceil n/m\rceil$ and the number of terms in \eqref{eq:anchor}
are the same for every alternative: a migration only moves the
coordinates, and the area changes only through the cell types.  If
$n_a\neq n_b$,
\begin{equation}
  \Delta A_c
  =\sum_{i\in C_{\sigma_c'}}W_iH_i-\sum_{i\in C_{\sigma_c}}W_iH_i
  \;\neq\;0,
  \qquad
  \text{footprint}
  =\Big\lceil\sqrt{n_a}\Big\rceil
   \times\Big\lceil n_a/\big\lceil\sqrt{n_a}\,\big\rceil\Big\rceil,
  \label{eq:size}
\end{equation}
and $\boldsymbol\alpha_c$ is a mean over $n_{\sigma_c'}$ cells instead of
$n_{\sigma_c}$.  The cost \eqref{eq:cost} contains no area term, so a
larger alternative can win its class on wirelength alone, and the
migration then grows the placed area without the selection having priced
the growth.
